无论如何，从 17 世纪末开始的微分学和积分学的历史分为两个部分。其中之一涉及这种演算的应用，它们总是越来越丰富、越来越众多、越来越多样。在平面曲线的微分几何、微分方程、幂级数、变分法（上文已经谈到过它们）之外，又加上了形状不规则的曲线继而曲面的微分几何、多重积分、偏微分方程、三角级数、众多特殊函数的研究，以及许多其他类型的问题。我们在这里只应关心那些就实变量函数而言有助于澄清、深化和巩固无穷小演算之基本原理的工作。

从这一观点看，18 世纪中叶的那些伟大论著没有提供多少新东西。英国的麦克劳林\(^{17}\){[}214{]}与大陆上的欧拉（{[}108 a{]}，(1)，第 X 至 XIII 卷）都忠实于各自所继承的传统。诚然，前者力图把牛顿的概念稍加澄清\(^{18}\)，而后者则把莱布尼茨形式主义推向极端，像莱布尼茨和泰勒一样，满足于把微分学建立在从差分演算出发的一种十分晦涩的过渡到极限之上，而对这种演算他恰好又给出了极为精致的阐述。但尤其重要的是，欧拉完成了莱布尼茨的事业，引入了 e、i 以及三角函数的那套沿用至今的记号并确保其被采纳，同时普及了记号 π。另一方面，尽管他大多数时候并不区分函数与解析表达式，但在谈到三角级数和弦振动问题时，他坚持认为，不能把自己局限于这样定义的函数（他把它们称为“连续的”），而是在需要时也应考虑由一条曲线的一段或若干段弧经验地给出的任意的、或者说“不连续的”函数（{[}108 a{]}，(1)，第 XXIII 卷，第 74–91 页）。最后，尽管这多少超出了我们的框架，在此仍不能不提到他把指数函数推广到复域，由此他得到了联系指数函数与三角函数的那些著名公式，以及复数对数的定义；由此最终阐明了对数与反圆函数之间那个著名的类比，或者用 17 世纪的语言说，圆的求积与双曲线的求积之间的类比——它早已被圣樊尚的格雷戈里观察到，被惠更斯、尤其被 J. 格雷戈里精确化，而在莱布尼茨和伯努利那里，它出现在对下式的形式积分之中

\[
\frac {1}{1 + x ^ {2}} = \frac {i}{2 (x + i)} - \frac {i}{2 (x - i)}.
\]

与此同时，达朗贝尔——无论在数学中还是在别处都是一切神秘主义的敌人——在一些出色的文章中（{[}75 a{]}，词条 DIFFÉRENTIEL 与 LIMITE，以及{[}75 b{]}）以最大的清晰性定义了极限和导数的概念，并坚决主张无穷小演算的“形而上学”归根结底尽在于此。但这一明智的忠告并未产生任何直接的效果。拉格朗日的煌煌巨著（{[}191{]}，第 IX–X 卷）代表了一种尝试，即把分析建立在牛顿诸概念中最可争议的一个之上，也就是把任意函数的概念与可展开为幂级数的函数的概念混为一谈的那个概念，并由此（通过考虑级数中一阶项的系数）引出微分的概念。当然，像拉格朗日这样水准的数学家在这种情况下不可能不获得重要而有用的成果，例如（以一种事实上独立于我们刚才所指出的出发点的方式）带积分形式余项的泰勒公式的一般证明，以及用中值定理对余项作出的估计；无论如何，拉格朗日的工作是魏尔斯特拉斯在复变函数论中方法的源头，也是现代形式级数代数理论的源头。但是，就其直接目标而言，它代表的与其说是进步，不如说是退步。

相反，有了柯西的那些教学著作（{[}56 a{]}，(2)，第 IV 卷），我们终于站到了坚实的地面上。他对函数的定义本质上与我们今天的一样，尽管所用的语言仍相当含糊。极限概念一经一劳永逸地确定，便被取为出发点；连续函数（取现代的意义）与导数的概念随即由此导出，连同它们的主要初等性质；而导数的存在性不再是一种信仰行为，而是变成用分析的寻常方法加以研究的对象。说实在的，柯西对此几乎不感兴趣；另一方面，虽然波尔查诺在独立达到同样一些原理之后构造了一个在任何一点都没有导数的连续函数的例子{[}27 a{]}，这个例子并未发表，问题直到 1872 年才由魏尔斯特拉斯在一部著作中公开地予以解决（{[}329 a{]}，第 II 卷，第 71–74 页）。

就积分而言，柯西的工作代表了一种向古代以及 17 世纪前半叶那些健全传统的回归，但所依靠的技术手段仍不充分。太久以来被降到第二位的定积分重新成为首要的概念，对这一概念，柯西使傅里叶所提出的记号 \(\int_{a}^{b} f(x) dx\) 得到了最终的采纳（以取代欧拉有时使用的那个笨拙记号 \(\int f(x) dx\left[\begin{array}{ccc}x & = & a \\ x & = & b\end{array}\right]\)）；而为了定义定积分，柯西回到了穷竭法，或者用我们今天的话说，回到“黎曼和”（把它称作阿基米德和或欧多克索斯和也许更好）。诚然，17 世纪认为没有必要对面积概念作一番批判性的检查，在他们看来这个概念至少与不可公度实数的概念一样清楚；但只要涉及的是单调曲线或逐段单调曲线，“黎曼”和向曲线下方面积的收敛，对 17 世纪所有关心严格性的作者——如费马、帕斯卡、巴罗——来说都是一个熟悉的思想；而 J. 格雷戈里由于对过渡到极限的思考以及已经十分抽象形式的“区间套”原理的熟悉而有了特别的准备，据推测他甚至拟定过一个细致的证明，它一直未曾发表（{[}136 d{]}，第 445–446 页），倘若柯西知道它，几乎可以不加任何修改地加以利用。\(^{19}\)不幸的是，柯西竟想对任意的连续函数证明积分的存在性，也就是说“黎曼和”的收敛性；而他的证明，倘若建立在闭区间上连续函数的一致连续性定理之上本会变成正确的，却因缺少这一概念而丧失了一切说服力。狄利克雷在撰写他关于三角级数的著名论文时似乎也没有看出这个困难，因为他在那里把上述定理引作“容易证明”的（{[}92{]}，第 I 卷，第 136 页）；确实，他最终只是把该定理应用于有界的逐段单调函数；黎曼更为审慎，只是在需要使用他关于“黎曼和”收敛的充分必要条件时才提到这类函数（{[}259 a{]}，第 227–271 页）。在海内建立了一致连续性定理（参见第 145 页）之后，这件事当然不再有任何困难；1875 年达布在他关于不连续函数积分的论文{[}77{]}中轻松地处理了它，在那篇论文中，他在许多地方与 P. 杜布瓦–雷蒙大约同时发表的重要研究不谋而合。同样地，连续函数之积分的线性性质第一次、而且这次是最终地得到了证明。另一方面，序列或级数一致收敛的概念（由塞德尔等人于 1848 年引入，并尤其由魏尔斯特拉斯所凸显（参见第 205 页））使人们得以在确实稍嫌过于狭窄的条件下，为级数的逐项积分和积分号 \(\int\) 下的求导建立一个坚实的基础，同时等待我们在这里无须谈论的那些现代理论，它们将以一种暂时成为定论的方式澄清这些问题。

这样我们就达到了古典无穷小演算的最后阶段，即 19 世纪末那些伟大的分析教程所代表的阶段；从我们所关心的观点看，若尔当的那部教程{[}174c{]}在其中占据着一个卓绝的地位，一方面是出于美学上的理由，但也因为它既是对古典分析结果的一个令人赞叹的铺陈，又在许多方面预告了现代分析并为它准备了道路。

\chapter{渐近展开.}

各阶“无穷小”（或“无穷大”）之间的区分，从关于微分学的最早著作、例如费马的著作起，就已经隐含地出现；它随着牛顿和莱布尼茨的“高阶差分”理论而变成完全自觉的；而且人们随即便观察到，在最简单的情形，形如 \(f(x)/g(x)\) 的表达式在 f 与 g 都趋于 0 的那种点处的极限（或“真值”），由这两个函数在所考虑点的邻域内的泰勒展开给出（“洛必达法则”，它很可能应归功于约翰·伯努利）。

撇开这个初等情形不谈，早在 17 世纪末就已经摆在数学家面前的“渐近展开”主要问题，是当 n 很大时形如 \(\sum_{k=1}^{n}f(k)\) 的和的精确或近似计算；这种计算无论对插值、对级数和的数值估计，还是在概率演算中，实际上都是必需的，在概率演算中，诸如 \(n!\) 或 \(\binom{a}{n}\) 这样的“大数的函数”起着举足轻重的作用。早在牛顿，为了在 n 大时得到 \(\sum_{k=1}^{n}\frac{1}{a+k}\) 的近似值，就指出了一种方法，它（在这个特殊情形）归结为计算欧拉–麦克劳林公式的前几项（{[}262{]}，第 II 卷，第 309–310 页）。临近该世纪末，雅各布·伯努利在他关于概率演算的研究中提出要确定和 \(S_{k}(n)=\sum_{p=1}^{n}p^{k}\)，即 n 的多项式\(^{1}\)，他发现了这些多项式的一般构成规律（但未给出证明），从而第一次在这些多项式系数的表达式中引入了以他的名字命名的那些数，以及使这些数得以计算的递推关系（{[}19 b{]}，第 97 页）。1730 年，斯特林通过一种按递推计算系数的程序，对 n 无限增长时的 \(\sum_{k=1}^{n}\log(x+ka)\) 得到了一个渐近展开。

欧拉关于级数及相关问题的决定性工作定年在 1730 至 1745 年之间。设 \(S(n) = \sum_{k=1}^{n} f(k)\)，他把泰勒公式应用于函数 \(S(n)\)，由此得到

\[
f (n) = S (n) - S (n - 1) = \frac {d S}{d n} - \frac {1}{2 !} \frac {d ^ {2} S}{d n ^ {2}} + \frac {1}{3 !} \frac {d ^ {3} S}{d n ^ {3}} - \dots ,
\]

他用待定系数法“反转”这个方程，同时寻找形如下式的解

\[
S (n) = \alpha \int f (n) d n + \beta f (n) + \gamma \frac {d f}{d n} + \delta \frac {d ^ {2} f}{d n ^ {2}} + \dots ;
\]

于是他逐步得到

\[
S (n) = \int f (n) d n + \frac {f (n)}{2} + \frac {1}{12} \frac {d f}{d n} - \frac {1}{720} \frac {d ^ {3} f}{d n ^ {3}} + \frac {1}{30,240} \frac {d ^ {5} f}{d n ^ {5}} - \dots
\]

起初他还不能确定系数的构成规律（{[}108 a{]}，(1)，第 XIV 卷，第 42–72 页及第 108–123 页）。但大约在 1735 年，通过与多项式分解为一次因式的类比，他毫不犹豫地写下了公式

\[
\begin{array}{c} 1 - \frac {\sin s}{\sin \alpha} = \\ \left(1 - \frac {s}{\alpha}\right) \left(1 - \frac {s}{\pi - \alpha}\right) \left(1 - \frac {s}{- \pi - \alpha}\right) \left(1 - \frac {s}{2 \pi + \alpha}\right) \left(1 - \frac {s}{- 2 \pi + \alpha}\right) \dots \end{array}
\]

令整个级数两个成员的展开式的系数相等，他便特别地（对 \(\alpha = \frac{\pi}{2}\)）得到了级数 \(\sum_{n=1}^{\infty} \frac{1}{n^{2k}}\) 的著名表达式\(^2\)，即用 \(\pi\) 的幂表出的那些（同上出处，第 73–86 页）。几年之后，他终于认识到这些 \(\pi\) 的幂的系数与他的求和公式的系数由同样的方程给出，并认出了它们与伯努利所引入的数、与 \(z/(e^z - 1)\) 的级数展开的系数之间的联系（同上出处，第 407–462 页）。

与欧拉相互独立，麦克劳林在大约同一时间、沿一条稍欠冒险而与今天所遵循的路径相近的途径，到达了同一个求和公式：他实际上是在迭代那个把 \(f(x)\) 用差 \(f^{(2k + 1)}(x + 1) - f^{(2k + 1)}(x)\) 表达出来的“泰勒式”公式，而这个公式他是用待定系数法“反转”这些差的泰勒展开而得到的（{[}214{]}，第 II 卷，第 672–675 页）；此外，他也没有注意到为欧拉所发现的系数构成规律。

但是，麦克劳林像欧拉以及他那个时代的所有数学家一样，把他所有的公式都作为级数展开来陈述，其收敛性甚至无人研究。这并不是说收敛级数的概念在当时完全受到忽视：由雅各布·伯努利已经知道调和级数发散，而欧拉本人在用他的求和公式估计这个级数前 n 项之和时也记录了这一结果（{[}108 a{]}，(1)，第 XIV 卷，第 87–100 页及第 108–123 页）；也是欧拉注意到两个相邻的伯努利数之比无限增长，由此可知以这些数为系数的整幂级数不可能收敛（同上出处，第 357 页）。\(^{3}\)但是形式计算的倾向更为强大，欧拉本人那非凡的直觉也不能阻止他有时堕入荒谬，例如当他写下 \(0 = \sum_{n=-\infty}^{+\infty} x^{n}\) 时（同上出处，第 362 页）。\(^{4}\)

我们在别处（见第 153 页）已经说过，19 世纪初的数学家们如何厌倦了这种无节制、无根据的形式主义，而把分析重新引回严格性的道路。一旦收敛级数的概念得到精确化，就出现了对一些简单判别法的需要，用以通过与已知的积分或级数相比较来证明积分和级数的收敛性；柯西在他的《代数分析》中给出了若干判别法（{[}56 a{]}，(2)，第 III 卷），而阿贝尔在一篇身后发表的论文中（{[}1{]}，第 II 卷，第 197–205 页）得到了对数判别法。另一方面，柯西（{[}56 a{]}，(1)，第 VIII 卷，第 18–25 页）澄清了诸如斯特林级数这类级数（由应用欧拉–麦克劳林公式得到，常被称为“半收敛级数”）的悖论：他证明，如果（由于欧拉关于伯努利数的评论）这样一个级数的通项 \(u_k(n)\) 在 \(n\) 取固定值时随 \(k\) 无限增长，那么尽管如此，对固定的 \(k\) 值，部分和 \(s_k(n) = \sum_{h=1}^{k} u_h(n)\) 仍给出级数所“代表”的函数的一个渐近展开（对 \(n\) 趋于 \(+\infty\)），而且 \(k\) 越大越精确。

在古典分析的大部分计算中，都有可能得到函数渐近展开的一般构成规律，其项数可以任意多；这一事实促成了级数与渐近展开之间一种持久的混淆（至少在语言上）；以致 H. 庞加莱在 1886 年费心把渐近展开的初等规则加以编纂时（{[}251 a{]}，第 I 卷，第 290–296 页）（依 \(1/x\) 的整幂、在 \(+\infty\) 的邻域中），使用的仍然是级数论的语言。只是随着来自解析数论的渐近展开的出现，才在渐近展开的概念与级数的概念之间实现了明确的区分，原因在于，在这门理论所处理的大多数问题中，所寻求的展开中只有很少几项（通常只有一项）能够显式地得到。

这些问题也使数学家们熟悉了变量（实变量或整数变量）的幂以外的比较尺度的使用。这一推广应追溯到 P. 杜布瓦–雷蒙的工作{[}94 a and b{]}，他第一个系统地处理了在一点的邻域中比较函数的问题，并在极具独创性的著作中认识到比较尺度的“非阿基米德”特征，同时以一般的方式研究了比较关系的积分与微分，并由此引出大量有趣的结论{[}94 b{]}。然而他的证明缺乏清晰性与严格性，杜布瓦–雷蒙这些结果的正确表述应归功于 G. H. 哈代{[}147{]}：他的主要贡献在于认识并证明了一类“初等函数”即函数 (H) 的存在性，对这类函数，分析的通常运算（尤其是微分）可以应用于比较关系。

\chapter{Γ 函数.}

对序列 \((u_{n})\) 用一个依赖于实参数 \(\lambda\) 且等于 \(u_{n}\)（当 \(\lambda = n\) 时）的积分的值来进行“插值”的想法，可追溯到沃利斯（参见第 187 页以下）。正是这个想法在 1730 年（{[}108 a{]}，(1)，第 XIV 卷，第 1–24 页）成为引导欧拉对阶乘序列进行插值的主要线索。他首先注意到 \(n!\) 等于无穷乘积 \(\prod_{k=1}^{\infty}\left(\frac{k+1}{k}\right)^{n}\frac{k}{k+n}\)，这个乘积对 n 的每个值（无论是否整数）都有定义，而且特别地，当 \(n = \frac{1}{2}\) 时由沃利斯公式它取值 \(\frac{1}{2}\sqrt{\pi}\)。这一结果与沃利斯诸结果的类比接着引导他重新拾起积分 \(\int_{0}^{1} x^{e}(1-x)^{n} dx\)（n 为整数，e 任意），它已被沃利斯考虑过。欧拉借助二项展开得到它的值 \(\frac{n!}{(e+1)(e+2)\cdots(e+n)}\)；随后一个变量代换向他表明，\(n!\) 是积分 \(\int_{0}^{1}\left(\frac{1-x^{s}}{z}\right)^{n} dx\) 当 z 趋于 0 时的极限，由此得到“第二类欧拉积分”\(n! = \int_{0}^{1}\left(\log\frac{1}{x}\right)^{n} dx\)；用同样的方法并利用沃利斯公式，他得到公式 \(\int_{0}^{1}\sqrt{\log\frac{1}{x}} dx = \frac{1}{2}\sqrt{\pi}\)。在其后的工作中，欧拉经常回到这些积分上来；他就这样发现了余元关系（{[}108 a{]}，(1)，第 XV 卷，第 82 页及第 XVII 卷，第 342 页）、公式 \(B(p,q) = \Gamma(p)\Gamma(q)/\Gamma(p+q)\)（{[}108 a{]}，(1)，第 XVII 卷，第 355 页）以及勒让德–高斯公式对应于 x = 1 的特例（{[}108 a{]}，(1)，第 XIX 卷，第 483 页）；当然，这一切都全然不关心收敛性的问题。

高斯在他关于超几何函数研究的背景下继续了对 \(\Gamma\) 函数的研究，\(\Gamma\) 函数是超几何函数的一个极限情形（{[}124 a{]}，第 III 卷，第 125–162 页）；正是在这项研究中他得到了一般的乘法公式（不久前勒让德已对 \(p = 2\) 记下了它）。其后关于 \(\Gamma\) 的工作主要涉及把这一函数推广到复域。直到最近人们才注意到，对数凸性在实域上把 \(\Gamma(x)\) 在函数方程 \(f(x + 1) = xf(x)\) 的全部解中确定到一个因子（{[}26{]}，第 149–164 页）；而阿廷已经表明{[}7 d{]}，关于 \(\Gamma(x)\) 的一切经典结果都可以简单地与这一性质联系起来。

\chapter{函数空间.}

我们知道，任意函数的概念大致直到 19 世纪初才出现。至于以一般的方式研究函数集合、并为它们配备拓扑结构的思想，则不早于黎曼（见第 141 页）出现，而且直到 19 世纪末才开始付诸使用，就更是如此了。

然而，数值函数序列收敛的概念，从无穷小演算的开端起就或多或少自觉地被使用着。但那只是简单收敛的问题，而且在波尔查诺和柯西精确地定义收敛级数和连续函数的概念之前，也不可能不是如此。后者最初并没有看出简单收敛与一致收敛之间的区别，还以为能够证明每一收敛的连续函数项级数都以连续函数为其和（{[}56 a{]}，(2)，第 III 卷，第 120 页）。这一错误几乎立即被阿贝尔揭露，他同时证明了每个整幂级数在其收敛区间的内部是连续的，所用的论证已成为经典，它在这一特殊情形本质上用到一致收敛的观念（{[}1{]}，第 I 卷，第 223–224 页）。剩下的只是以一般的方式把这一观念分离出来，这由斯托克斯和塞德尔在 1847–48 年、以及柯西本人在 1853 年（{[}56 a{]}，(1)，第 XII 卷，第 30 页）独立地完成。\(^{1}\)

在魏尔斯特拉斯和黎曼的影响下，一致收敛概念及相关问题的系统研究在 19 世纪最后三分之一的时间里由德国学派（汉克尔、杜布瓦–雷蒙）尤其是意大利学派发展起来：迪尼和阿尔泽拉明确了连续函数序列的极限为连续的必要条件，而阿斯科利则引入了等度连续性这一基本概念，并证明了刻画连续函数之紧集的定理{[}10{]}（这个定理后来由蒙泰尔在他的“正规族”理论中加以普及，正规族无非就是解析函数的相对紧集）。

另一方面，魏尔斯特拉斯本人（{[}329 a{]}，第 III 卷，第 1–37 页）发现了用多项式在一个有界集上一致逼近一个或多个实变量的连续数值函数的可能性，这一结果立即激起了强烈的兴趣，并引出了许多“定量”研究，它们处于我们目前所遵循的观点之外。

这些问题上的现代贡献，首先在于赋予它们以所能有的全部广度，即把它们用于定义集和值集不再局限于 R 或有限维空间的函数，并借助一般的拓扑概念把它们安置到其自然的框架之中。特别地，已在古典分析中显示出是一阶工具的魏尔斯特拉斯定理，近年来被 M. H. 斯通推广到了远为一般的情形：通过发展 H. 勒贝格引入的一个想法（在魏尔斯特拉斯定理的一个证明中），他突显了格在连续数值函数逼近中所起的重要作用（用“格多项式”逼近），另一方面又表明魏尔斯特拉斯定理的推广形式如何立即蕴涵一整串类似的逼近定理，它们就这样以一种连贯得多的方式被归拢在一起{[}301 c{]}。

\chapter{拓扑向量空间.}

拓扑向量空间的一般理论大约创立于 1920 至 1930 年间。但在此之前，它早已由泛函分析中众多问题的研究作了长期的准备；而且，若不至少概括地指出这些问题的研究如何使数学家们（尤其从 20 世纪初起）逐渐认识到所论问题的共同渊源，以及把它们表述得远为一般、并对它们应用统一的求解程序的可能性，就无法追溯这段历史。

可以说，代数与分析之间的类比，以及把函数方程（也就是说未知量是一个函数的方程）看作代数方程的“极限情形”的想法，都追溯到无穷小演算的开端，这在某种方式上回应了这种“从有限到无限”的推广需要。但无穷小演算在代数上的直接祖先是有限差分演算（参见第 185 页以下），而不是一般线性方程组的求解；直到 18 世纪中叶，后者与微分学诸问题之间的最初类比才在弦振动方程的背景下显现出来。我们在这里不打算深入这个问题的历史细节；但我们必须指出两个基本思想的出现，它们在下文中还会不断出现，而且两者似乎都应归功于 D. 伯努利。第一个思想在于把弦的振动看作 n 个质点组成的系统的振动当 n 趋于无穷时的“极限情形”；众所周知，当 n 有限时，这个问题在稍后给出了寻求线性变换特征值的第一个例子（参见第 88 页）；在所考虑的“过渡到极限”中，与这些数相对应的是弦的“固有振动”的频率，它们早在此前很久就已被实验观察到，而其理论上的存在性已在本世纪初（特别是由泰勒）建立。这个形式上的类比，尽管在后来的文献中相当少被提及（{[}302 b{]}，第 390 页），在 19 世纪期间似乎从未被忽视；但正如我们将在后文看到的，直到大约 1890–1900 年它才获得其全部重要性。

D. 伯努利的另一个想法（也许是受实验事实的启发）是“叠加原理”，按照这个原理，弦的最一般的振动必须能够“分解”为“固有振动”的叠加，用数学的语言说，这意味着弦振动方程的一般解必须能够展开为级数 \(\sum_{n} c_{n} \phi_{n}(x, t)\)，其中 \(\phi_{n}(x, t)\) 表示各个固有振动。众所周知，这一原理后来引发了一场关于把“任意”函数展开为三角级数之可能性的长期论战，这场论战直到 19 世纪前三分之一的时间里才由傅里叶和狄利克雷的工作了结。但甚至在这一结果达到之前，人们已经遇到过“正交”\(^{1}\)函数展开为级数的其他例子：球面函数和勒让德多项式，以及形如 \((e^{i\lambda_{n}x})\) 的各种系统，其中 \(\lambda_{n}\) 不再是同一个数的倍数，这些系统从 18 世纪起就被引入振动问题之中，也被傅里叶和泊松在关于热理论的研究中引入。大约在 1830 年，斯图姆{[}302 a and b{]}和刘维尔{[}204 a and b{]}把这些各别特殊情形中观察到的全部现象系统化为一个关于单变量函数的一般振动理论：他们考虑微分方程

\[
\frac {d}{d x} \left(p (x) \frac {d y}{d x}\right) + \lambda \rho (x) y = 0 (p (x) > 0, \rho (x) > 0)\tag{21.1}
\]

连同边界条件

\[
\left\{ \begin{array}{l l} y ^ {\prime} (a) - h _ {1} y (a) & = 0 \\ y ^ {\prime} (b) + h _ {2} y (b) & = 0 \end{array} \right. (h _ {1} \geq 0, h _ {2} \geq 0, a <   b)\tag{21.2}
\]

并证明以下基本结果：

\begin{enumerate}
\def\labelenumi{\arabic{enumi})}
\item
  这个问题具有 \(\neq 0\) 的解，仅当 \(\lambda\) 取到某个数列 \((\lambda_{n})\)（其中的数 \textgreater{} 0 且趋于 \(+\infty\)）中的值；
\item
  对每个 \(\lambda_{n}\)，解都是同一个函数 \(v_{n}\) 的倍数，这个函数可以假定已由条件 \(\int_{a}^{b} \rho v_{n}^{2} dx = 1\) 加以“规范化”，而且我们有 \(\int_{a}^{b} \rho v_{m} v_{n} dx = 0\)（对 \(m \neq n\)）；
\item
  每个函数 \(f\)，若它在 \([a, b]\) 中二次可微且满足边界条件 (21.2)，则可展开为一致收敛的级数 \(f(x) = \sum_{n} c_{n} v_{n}(x)\)，其中 \(c_{n} = \int_{a}^{b} \rho f v_{n} dx\)；
\item
  我们有等式 \(\int_{a}^{b}\rho f^{2}dx=\sum_{n}c_{n}^{2}\)（帕塞瓦尔已于 1799 年——而且是以纯粹形式的方式——对三角函数系证明过它，由它立即可得出后者（仍就三角级数）于 1828 年陈述的“贝塞尔不等式”）。
\end{enumerate}

半个世纪之后，这些性质由格拉姆的工作{[}133{]}加以补全；格拉姆继切比雪夫的研究之后，揭示了正交函数级数展开与“最佳平方逼近”问题（它直接源出于高斯在误差理论中的“最小二乘法”）之间的关系：给定函数的一个有限序列 \((\psi_{i})_{1\leq i\leq n}\)，这个问题的内容是，对函数 f 求线性组合 \(\sum_{i}a_{i}\psi_{i}\)，使积分 \(\int_{a}^{b}\rho(f-\sum_{i}a_{i}\psi_{i})^{2}dx\) 达到最小。原则上这只是一个基本的线性代数问题，但格拉姆以独创的方式解决了它，即对 \(\psi_{i}\) 施行通常以埃哈德·施密特之名著称的“规范正交化”程序。接下来转到无穷正交系 \((\phi_{n})\) 的情形，他给自己提出的问题是：用序列前 n 个函数的线性组合对函数 f 所作的“最佳平方逼近” \(\mu_{n}\)，何时当 n 趋于无穷时趋于 0；\(^{2}\)他由此被引导去定义完备规范正交系的概念，并认识到这一性质等价于不存在 \(\neq0\) 的、与所有 \(\phi_{n}\) 都正交的函数。他甚至试图阐明“均方收敛”的概念，但是，在测度论的基本概念被引入之前，他在这个方向上很难得到任何十分明确的结果。

在 19 世纪下半叶，分析学家们的主要力量主要用于把斯图姆–刘维尔理论推广到多变量函数，数理物理中椭圆型偏微分方程的研究以及自然伴随它们的边值问题尤其引导了这一推广。兴趣主要集中在“振动膜”方程

\[
L _ {\lambda} (u) \equiv \Delta u + \lambda u = 0\tag{21.3}
\]

其中所寻求的是在相当正规的区域 \(G\) 的边界上为零的解；这个问题所呈现的巨大分析困难只是一点一点地被克服的，人们原本认为不可能把对单变量函数行之有效的方法应用于这个问题。让我们回顾通向解决的主要阶段：\(G\) 的“格林函数”的引入，其存在性由施瓦茨证明；最小特征值的存在性的证明，同样出自施瓦茨；最后，1894 年，在一篇著名的论文中（{[}251 a{]}，第 IX 卷，第 123–196 页），H. 庞加莱成功地证明了全部特征值的存在性及其本质性质，其办法是：对给定的“右端” \(f\)，考虑这样的解 \(u_{\lambda}\)：它满足方程 \(L_{\lambda}(u) = f\) 且在边界上为零，并通过巧妙地推广施瓦茨的方法证明 \(u_{\lambda}\) 是复变量 \(\lambda\) 的亚纯函数，只有实的单极点 \(\lambda_n\)，它们正是所寻求的特征值。

这些研究与线性积分方程论的开端紧密相连，后者对现代观念之到来所作的贡献无疑最大。我们在这里只限于对这门理论的发展给出一些简短的说明。这类函数方程最初在 19 世纪上半叶零星地出现（阿贝尔、刘维尔），自从比尔和 C. 诺伊曼把相当正规的区域 G 上“狄利克雷问题”的求解化为一个“第二类积分方程”

\[
u (x) + \int_ {a} ^ {b} K (x, y) u (y) d y = f (x)\tag{21.4}
\]

的求解（未知函数为 u）；这个方程 C. 诺伊曼已于 1877 年借助“逐次逼近”程序解出。无疑既受到已提到的代数类比的驱使，也受到他刚刚在振动膜方程上所得结果的驱使，H. 庞加莱在 1896 年（{[}251 a{]}，第 IX 卷，第 202–272 页）想到在前述方程的积分之前引入一个变参数 \(\lambda\)，并断言，如同在振动膜方程的情形一样，解这时是 \(\lambda\) 的亚纯函数；但他没有能证明这一结果，它（对连续的“核” K 和有限区间 \([a, b]\)）直到四年后才由 I. 弗雷德霍姆建立{[}116{]}。后者也许甚至比他的前辈们更为自觉，让自己完全受 (21.4) 与线性方程组

\[
\sum_ {q = 1} ^ {n} \left(\delta_ {p q} + \frac {1}{n} a _ {p q}\right) x _ {q} = b _ {p} (1 \leq p \leq n)\tag{21.5}
\]

的类比的引导，以便把 (21.4) 的解表示成两个表达式之商，这两个表达式是仿照克拉默公式中出现的行列式构造出来的。这也不是一个新想法：早在 19 世纪初，“待定系数法”（其内容是：对一个假定可展开为级数 \(\sum_{n} c_{n} \phi_{n}\) 的未知函数——其中 \(\phi_{n}\) 是已知函数——通过计算系数 \(c_{n}\) 来求得它）就已引出了“无穷多个未知数的线性方程组”

\[
\sum_ {j = 1} ^ {\infty} a _ {i j} x _ {j} = b _ {i} (i = 1, 2, \dots).\tag{21.6}
\]

傅里叶碰到这样一个方程组时，仍像 18 世纪的数学家那样去“解”它：他把指标 i 或 j 大于 n 的项统统删去，用克拉默法则显式地解出这样得到的有限方程组，然后“过渡到极限”，在解中令 n 趋于 +∞！到后来，当这种偷懒的把戏不再够用时，对这一问题的最初尝试仍然是通过行列式理论作出的；从 1886 年起（继希尔的工作之后），H. 庞加莱、然后是 H. 冯·科赫，建立了一种“无穷行列式”理论，使某些类型的方程组 (21.6) 可以按古典的模式求解；即使这些结果不能直接适用于弗雷德霍姆所瞄准的问题，至少可以肯定，冯·科赫的理论特别曾作为构造他的“行列式”的范本。

正是在这个时刻希尔伯特登场，给这门理论以新的推动{[}163 b{]}。他首先完成弗雷德霍姆的工作，有效地实现了从 (21.5) 的解到 (21.4) 的解的过渡到极限；但他随即为实二次型的理论补上了相应的过渡到极限，具有对称核（也就是说 \(K(y, x) = K(x, y)\)）的那类积分方程自然导向这个理论，而且这类方程在数理物理中最为常见。他就这样得到了那个直接推广二次型约化到主轴的基本公式

\[
\int_ {a} ^ {b} \int_ {a} ^ {b} K (s, t) x (s) x (t) d s d t = \sum_ {n = 1} ^ {\infty} \frac {1}{\lambda_ {n}} \left(\int_ {a} ^ {b} \phi_ {n} (s) x (s) d s\right) ^ {2},\tag{21.7}
\]

其中 \(\lambda_{n}\) 是核 \(K\) 的（必然为实的）特征值，\(\phi_{n}\) 构成相应本征函数的正交系，而且当 \(\int_{a}^{b} x^{2}(s) dx \leq 1\) 时公式 (21.7) 的右端是一个收敛级数。他还表明，每个可“表示”为形式 \(f(x) = \int_{a}^{b} K(x, y) g(y) dy\) 的函数如何容许“展开” \(\sum_{n=1}^{\infty} \phi_{n} \int_{a}^{b} \phi_{n}(y) f(y) dy\)，而且循着与二次型经典理论的类比，他指出了一个用变分方法确定 \(\lambda_{n}\) 的程序，它无非是二次曲面主轴那个著名的极值性质的推广（{[}163 b{]}，第 1–38 页）。

希尔伯特的这些最初结果几乎立即被 E. 施密特以更简单也更一般的形式接手，避免了“弗雷德霍姆行列式”的引入以及从有限到无限的过渡，而且已经非常接近抽象的陈述，积分的线性性与正性这两个基本性质显然只在证明中被使用{[}274 a{]}。但希尔伯特那时已经达到了还要更为一般得多的观念。此前的一切工作凸显了平方可积函数的重要性，而帕塞瓦尔公式在这些函数与序列 \((c_{n})\)（它们满足 \(\sum_{n} c_{n}^{2} < +\infty\)）之间建立了紧密的联系。无疑正是这个想法引导着希尔伯特在 1906 年的论文（{[}163 b{]}，第 XI–XIII 章）中重新拾起古老的“待定系数法”，证明积分方程 (21.4) 的求解等价于一个由无穷多个线性方程组成的方程组

\[
x _ {p} + \sum_ {q = 1} ^ {\infty} k _ {p q} x _ {q} = b _ {p} (p = 1, 2, \dots)\tag{21.8}
\]

其中未知函数 u 的“傅里叶系数”为 \(x_{p} = \int_{a}^{b} u(t) \omega_{p}(t) dt\)（相对于一个给定的完备规范正交系 \((\omega_{n})\)）（其中 \(b_{p} = \int_{a}^{b} f(t) \omega_{p}(t) dt\)，\(k_{pq} = \int_{a}^{b} \int_{a}^{b} K(s, t) \omega_{p}(s) \omega_{q}(t) ds dt\)。而且，从这个观点看，(21.8) 中唯一值得考虑的解是那些使 \(\sum_{n} x_{n}^{2} < +\infty\) 的解；希尔伯特也系统地把自身局限于这种类型的解；但与此相对，他放宽了加于“无穷矩阵”（\(k_{pq}\)）上的条件（在 (21.8) 中该矩阵满足 \(\sum_{p,q} k_{pq}^{2} < +\infty\)）。从那时起，事情很清楚：由实数序列 \(x = (x_{n})\)（满足 \(\sum_{n} x_{n}^{2} < +\infty\)）组成的“希尔伯特空间”虽然未被明确引入，却潜藏于整个理论之下，并且显示为从有限维欧氏空间出发的一个“过渡到极限”。除此以外，对后继的发展特别重要的是，希尔伯特被引导在这个空间中引入的不仅是一种、而是两种不同的收敛概念（相应于所谓弱拓扑和强拓扑\(^{3}\)），以及一个“选择原理”，它无非是单位球的弱紧性。他在求解方程组 (21.8) 的背景下发展的新线性代数，完全建立在这些拓扑概念之上：线性映射、线性型和双线性型（与线性映射相关联）按照它们的“连续性”\(^{4}\)性质被分类和研究。特别地，希尔伯特发现弗雷德霍姆方法的成功依赖于“全连续”概念，他通过就双线性型对它进行表述\(^{5}\)并作深入研究来揭示这一概念；我们在这里无法就这一重要概念给出更多细节，也无法详述希尔伯特在其间开创对称双线性型（有界或否）谱论的那令人赞叹的深刻工作。

希尔伯特的语言仍是古典的，而且在整个《线性积分方程一般理论》中，他始终注视着理论的应用，并发展了众多的例子（约占全卷一半）。下一代人已经采取了更为抽象的观点。在弗雷歇和 F. 里斯关于一般拓扑的思想影响下（见第 142 页），E. 施密特{[}274 b{]}和弗雷歇本人在 1907–1908 年间有意把欧氏几何的语言引入“希尔伯特空间”（实的或复的）；正是在这些著作中第一次提到范数（用现在的记号 \(||x||\)）、它所满足的三角不等式、希尔伯特空间“可分”且完备的事实；此外，E. 施密特证明了到闭线性簇上的正交投影的存在性，这使他能把希尔伯特的线性方程组理论表述为更简单也更一般的形式。同样在 1907 年，弗雷歇和 F. 里斯注意到平方可和序列的空间有一个完全类似的“几何”；几个月后，当 F. 里斯和 E. 菲舍尔证明这个空间完备且同构于“希尔伯特空间”时，这一类比得到了完美的解释，同时也以令人惊异的方式凸显了勒贝格新创造的工具的价值。从那时起，希尔伯特空间理论的要点即可视为已被获得；在较晚近的进展中，最值得一提的是 M. H. 斯通和 J. 冯·诺伊曼大约在 1930 年给出的理论的公理化表述，以及大约 1934 年在雷利希、勒维希和 F. 里斯的工作中对“可分性”限制的摆脱{[}260 g{]}。

然而，在 20 世纪的最初几年，另外一些思想潮流也在涌现，以加强导向赋范空间理论的趋势。“泛函”的一般观念（也就是定义在一个集合上的取数值的函数，该集合的元素本身是一个或多个实变量的数值函数）出现于 19 世纪最后几十年，一方面与变分法相联系，另一方面与积分方程论相联系。如果说这一观念以及更一般的“算子”观念的揭示主要应归功于以平凯勒、尤其是沃尔泰拉为中心的意大利学派，那么这个学派的工作往往仍带有相当形式的性质，并且局限于特殊类型的问题，缺乏对底层拓扑概念足够深入的分析。1903 年，阿达马开创了“拓扑”对偶的现代理论，他寻求紧区间 I 上数值函数的空间 \(\mathcal{C}(I)\)（配备一致收敛拓扑的空间）上最一般的连续线性“泛函”，并把它们刻画为积分序列 \(x \mapsto \int_{I} k_{n}(t)x(t)dt\) 的极限。1907 年，弗雷歇和 F. 里斯以同样的方式表明，在希尔伯特空间上，连续线性型就是希尔伯特引入的“有界”型；随后在 1909 年，F. 里斯把阿达马定理定型，用斯蒂尔杰斯积分表达了 \(\mathcal{C}(I)\) 上的一切连续线性泛函，这个定理后来成为现代积分论的出发点（见第 226 页）。

第二年，还是 F. 里斯{[}260 c{]}在理论上作出了新的重要进展，他引入并研究了（仿照希尔伯特空间理论）区间 I 上 p 次可和函数的空间 \(L^{p}(I)\)（指数 p 满足 \(1 < p < +\infty\)），三年后他继之以{[}260 e{]}关于序列空间 \(L^{p}(\mathbf{N})\) 的类似工作；正如我们将在后文看到的，这些研究大大有助于澄清关于对偶的观念，因为在这里第一次遇到了两个互相对偶却并非自然同构的空间。\(^{6}\)

从那时起，F. 里斯就在考虑一个涵盖所有这些结果的公理化研究（{[}260 c{]}，第 452 页）；而且似乎只是出于一位分析学家不愿过分背离古典数学的顾虑，才使他没有以这种形式写出他 1918 年那篇关于弗雷德霍姆理论的著名论文{[}260 f{]}。他在那里原则上考虑紧区间上连续函数的空间 \(\mathcal{C}(I)\)；但在定义了这个空间的范数、并注意到 \(\mathcal{C}(I)\) 连同这个范数是完备的之后，他在论证中就再没有使用过完备赋范空间公理以外的任何东西。\(^{7}\)在这里不对这项工作作详细的考察，只须指出：正是在那里第一次以一般的方式定义了全连续线性映射的概念（通过把一个邻域变为相对紧集的性质）；\(^{8}\)而且，凭借公理化分析的一项杰作，弗雷德霍姆理论的全部（就其定性方面而言）被化归为单一的基本定理，即每个局部紧的赋范空间都是有限维的。

赋范空间的一般定义是 1920–22 年间由 S. 巴拿赫、H. 哈恩和 E. 黑利给出的（后者只考虑实数或复数序列的空间）。在其后的十年里，空间理论主要围绕两个对应用具有根本重要性的问题发展：对偶理论以及与贝尔“纲”概念相联系的诸定理。

我们已经看到，（拓扑意义上的）对偶观念可追溯到 20 世纪初；它潜存于希尔伯特的理论之下，并在 F. 里斯的工作中占据中心位置。例如后者早在 1911 年（{[}260 d{]}，第 41–42 页）就注意到，关系 \(|f(x)| \leq M||x||\)（在希尔伯特空间中被取作“有界”线性泛函的定义）当我们在空间 \(\mathcal{C}(I)\) 中时等价于 f 的连续性，而且所用的是一个具有完全一般性质的论证。在关于 \(\mathcal{C}(I)\) 上连续线性泛函的刻画方面，他还注意到，集合 A 在 \(\mathcal{C}(I)\) 中稠密的条件是：不存在 I 上与 A 的每个函数都“正交”的斯蒂尔杰斯测度 \(\mu \neq 0\)（从而推广了完备规范正交系的格拉姆条件）；最后，在同一项工作中他指出，空间 \(L^{\infty}\) 的对偶比斯蒂尔杰斯测度的空间“更大”（{[}260 d{]}，第 62 页）。

另一方面，在他关于空间 \(L^{p}(I)\) 和 \(L^{p}(\mathbf{N})\) 的工作中，F. 里斯成功地对 E. 施密特{[}274 b{]}给出的求解希尔伯特空间中线性方程组的方法作了修改，使之适用于更一般的情形。E. 施密特的想法在于，通过在由方程组 (21.6) 所表示的闭线性簇中寻求到原点的距离最小的点，来确定 (21.6) 的一个“极值”解。利用同一想法，F. 里斯证明，存在函数 \(x \in L^{p}(a, b)\) 满足方程组

\[
\int_ {a} ^ {b} \alpha_ {i} (t) x (t) d t = b _ {i} (i = 1, 2, \dots)\tag{21.9}
\]

（其中 \(\alpha_{i}\) 属于 \(L^{q}\)（满足 \(\frac{1}{p} +\frac{1}{q} = 1\)），而且还满足 \(\int_{a}^{b}|x(t)|^{p}dt\leq M^{p}\)），其必要充分条件是：对实数的所有有限序列 \((\lambda_i)_{1\leq i\leq n}\)，有

\[
\left| \sum_ {i = 1} ^ {n} \lambda_ {i} b _ {i} \right| \leq M. \left(\int_ {a} ^ {b} \sum_ {i = 1} ^ {n} | \lambda_ {i} \alpha_ {i} (t) | ^ {q} d t\right) ^ {1 / q}.\tag{21.10}
\]

1911 年{[}260 d{]}，他以类似的方式处理“广义矩问题”，即求解方程组

\[
\int_ {a} ^ {b} \alpha_ {i} (t) d \xi (t) = b _ {i} (i = 1, 2, \dots)\tag{21.11}
\]

其中 \(\alpha_{i}\) 连续而未知量是一个斯蒂尔杰斯测度 \(\xi;^{9}\)；这里可以清楚地看到，这个问题可以这样来陈述：按一个连续线性泛函在空间 \(\mathcal{C}(I)\) 中一列给定点上的值来确定这个泛函。黑利在 1912 年处理的也是这个形式的问题——他以一种颇为不同、适用范围更广的方法得到了 F. 里斯的条件\(^{10}\)——并在 1921 年以远为一般的条件重新处理了它。如我们上文所见，他引入了（序列空间上的）范数概念，并注意到这个概念推广了 n 维空间中凸体的“度规”，它被闵可夫斯基用于他关于“数的几何”的著名著作{[}221 a and b{]}。在这项工作中，闵可夫斯基还（在 \(R^{n}\) 中）定义了支撑超平面和“支撑函数”的概念{[}221 b{]}，并证明了在凸体的所有边界点处支撑超平面的存在性（{[}221 a{]}，第 33–35 页）。黑利把这些概念推广到配备任意范数的一个序列空间 E；他在 E 与空间 \(E'\)（由满足下列条件的序列 \(u = (u_n)\) 组成）之间建立了一种对偶：对每个 \(x = (x_n) \in E\)，级数 \((u_n x_n)\) 应当收敛；以 \(\langle u, x \rangle\) 记这个级数的和，他在 \(E'\) 中用公式 \(\sup_{x \neq 0} |\langle u, x \rangle| / ||x||\) 定义了一个范数，它在有限维空间中给出的正是支撑函数。\(^{11}\)E 中方程组 (21.6) 的解——其中每个序列 \(u_i = (a_{ij})_{j \geq 1}\) 都假定属于 \(E'\)——如黑利当时所证明的，化归为相继求解下列两个问题：\(1^\circ\) 在赋范空间 \(E'\) 上找一个连续线性型 L，使对每个指标 i 有 \(L(u_i) = b_i\)，如他所指出的，这导致 (21.10) 型的条件；\(2^\circ\) 寻求这样一个线性型能否写成 \(u \mapsto \langle u, x \rangle\)（对某个 \(x \in E\)）。正如黑利所观察到的，最后这个问题即使当 L 存在时也未必有解，他只限于给出一些在某些特殊情形蕴含解 \(x \in E\) 存在的充分条件{[}156{]}。

这些观念于 1927 年在 H. 哈恩的一篇基本论文{[}142{]}中获得了它们最终的形式，其结果两年后被 S. 巴拿赫（独立地）重新得到{[}14 c{]}。哈恩把闵可夫斯基–黑利程序应用于任意的赋范空间，从而赋予对偶空间一个（完备）赋范空间的结构；这使哈恩能立即考虑赋范空间的逐次对偶，并以一般的方式陈述黑利已预见到的自反空间问题。但尤其是，在保持范数的条件下延拓连续线性泛函这个压轴问题，由哈恩以完全一般的方式最终解决，用的是对维数作超限归纳的论证——由此给出了选择公理应用于泛函分析的一个最早的重要例子。\(^{12}\)在这些结果之外，巴拿赫又对连续线性映射与其转置之间的关系作了透彻的研究，借助一个关于对偶空间中弱闭子集的深刻定理，把直到那时只对空间 \(L^{p}\){[}260 c{]}已知的结果推广到一般的赋范空间；这些结果若使用几年后由豪斯多夫和巴拿赫本人引入的赋范空间的商空间概念来表述，就更为醒目。最后，又是巴拿赫发现了单位球的弱紧性（如我们上文所指出的，在许多特殊情形都被观察到）与自反性之间的联系，至少对可数型的空间是如此（{[}14 a{]}，第 189 页）。从此刻起，赋范空间的对偶理论就可以认为在其主要方面已经定型。

与此同时，一些外表悖谬的定理得到了澄清，其最早的例子可追溯到 1910 年前后。事实上，黑林格和托普利茨在那一年实质上证明了：希尔伯特空间上有界双线性型的一个序列 \(B_n(x,y)\)，若其值 \(B_n(a,b)\) 对每个给定的对 \((a,b)\) 都有界（界数先验地依赖于 \(a\) 和 \(b\)），则它在每个球上事实上都一致有界。他们的证明用反证法，其内容是用一种后来以“滑动驼峰法”著称的归纳法构造出一个违反假设的特殊对 \((a,b)\)，这个方法至今仍用于许多类似的问题。早在 1905 年，勒贝格就曾用类似的程序证明存在其傅里叶级数在某些点发散的连续函数；而且，在与黑林格和托普利茨同一年，他用同一方法证明了 \(L^1\) 中弱收敛的级数按范数有界。\(^{13}\)这类例子在随后的年代里成倍增加，但直到 1927 年都没有引入任何新观念；这一年，巴拿赫和施坦豪斯（有 S. 萨克斯的部分合作）把这些现象与贫集的概念以及完备度量空间上的贝尔定理联系起来，得到了一个涵盖此前所有特殊情形的一般命题{[}15{]}。此外，对完备赋范空间中“纲”问题的研究还同时把巴拿赫引向关于连续线性映射的许多其他结果；最引人注目也无疑最深刻的是“闭图象”定理，它与巴拿赫–施坦豪斯定理一样，在现代泛函分析中显示出是一阶工具{[}14 c{]}。

巴拿赫关于《线性算子》的论著的出版{[}14 a{]}，可以说标志着赋范空间理论成年期的开始。我们一直在谈论的所有结果以及许多其他结果都铺陈于这部书中，方式仍有些杂乱，但伴有取自分析各个领域的许多醒目的例子，它们似乎预示着这门理论辉煌的前景。事实上，这部著作取得了巨大的成功，其最直接的效果之一就是巴拿赫所用的语言和记号几乎被普遍采用。但是，尽管四十年来对巴拿赫空间进行了大量的研究（{[}203 bis{]}），如果不把巴拿赫代数理论及其在交换和非交换调和分析中的应用算作例外，那么这门理论对古典分析重大问题的几乎完全缺失的新应用，多少使人们对它寄予的希望落空了。

最为富有成效的发展，毋宁说是沿着对与赋范空间有关的概念作一种日益扩张、日益深入的公理化分析的方向发生的。虽然自 20 世纪初以来遇到的功能空间大都天生带有“自然的”范数，但也不是没有引起人们注意到一些例外。大约在 1910 年，E. H. 摩尔曾提议推广一致收敛的概念，即以“相对一致收敛”的概念来代替它，其中 0 的一个邻域由满足关系 \(|f(t)| \leq \epsilon g(t)\) 的函数 f 组成，g 是一个处处 \textgreater{} 0 的函数，它可以随邻域而变化。另一方面，在 1930 年之前人们已经注意到，像简单收敛、可测函数的依测度收敛或整函数的紧收敛这样的概念，都不能借助范数来定义；而且在 1926 年，弗雷歇已注意到这种性质的向量空间可以是可度量化的和完备的。但这些更一般空间的理论，只有与凸性思想相结合才能富有成效地发展。后者（我们在黑利那里已看到它的出现）曾是巴拿赫及其学生们研究的对象，他们认识到有可能以此把赋范空间理论中的许多命题解释得更具几何意味，从而为 J. 冯·诺伊曼 1935 年给出的局部凸空间的一般定义准备了道路。这些空间的理论，特别是涉及对偶的那些问题，最近十年间得到了特别的发展。在这方面必须指出的，一方面是一般拓扑学基本概念的建立（在 1930 至 1940 年间完成）所带来的简单性与一般性方面的进步；另一方面是有界集概念所取得的重要性，它由科尔莫哥洛夫和冯·诺伊曼于 1935 年引入，其在的对偶理论中的根本作用已由麦基{[}213 a and b{]}和格罗滕迪克{[}138 b and c{]}的工作所揭示。最后尤其是，可以肯定，激励这项研究的主要动力来自把理论应用于分析的新可能性，即巴拿赫理论无能为力的那些领域：这里必须提到由克特、托普利茨及其学生们自 1934 年起在一系列论文中发展的序列空间理论{[}184{]}，范塔皮耶“解析泛函”理论的最近建立，尤其是 L. 施瓦茨的广义函数理论{[}280{]}，现代局部凸空间理论在这里找到了一个无疑还远未穷竭的应用领域。

\chapter{局部紧空间上的积分.}

现代积分概念的发展，与函数观念的演进以及自 19 世纪初以来对实变量数值函数的深入研究紧密相连。众所周知，欧拉对函数概念的构想已经相当一般，因为在他看来，一条被每条平行于 Oy 轴的直线只交于一点的给定“任意”曲线就定义了一个函数 \(y = f(x)\)（参见第 196 页）；但是，与其同时代的大多数人一样，他拒绝承认这样的函数可以“解析地”表出。直到傅里叶的工作之前，这一观点几乎没有什么变化；而傅里叶所发现的把不连续函数表示为三角级数之和的可能性\(^{1}\)，对后世产生了决定性的影响。说实在的，傅里叶的证明完全缺乏严格性，其有效范围也未显现清楚；然而积分公式

\[
a _ {n} = \frac {1}{\pi} \int_ {- \pi} ^ {+ \pi} \phi (x) \cos n x d x, b _ {n} = \frac {1}{\pi} \int_ {- \pi} ^ {+ \pi} \phi (x) \sin n x d x (n \geq 1)\tag{22.1}
\]

它们给出函数 \(\phi\) 展开为傅里叶级数时的系数；只要假定 \(\phi\) 连续且逐段单调，这些公式便有一种直观上显而易见的意义。\(^{2}\)因此狄利克雷起初正是把自己局限于这些函数，在他建立傅里叶级数收敛性的那篇著名论文（{[}92{]}，第 I 卷，第 117–132 页）中即是如此；但在这项工作的末尾，他已经关心把他的结果推广到更广的函数类去了。众所周知，正是在这个场合，狄利克雷把傅里叶的观念精确化，定义了我们今天所理解的函数的一般概念；首先要澄清的一点自然是要知道，在哪些情形下仍能赋予公式 (22.1) 以意义。“当连续性{[}关于 \(\phi\) 的{]}的解有无穷多个\ldots{}”，狄利克雷说（同上出处，第 131–132 页），“那时就必须要求函数 \(\phi(x)\) 具有这样的性质：若以 \(a\) 和 \(b\) 表示 \(-\pi\) 与 \(+\pi\) 之间的两个任意量，则总能在 \(a\) 与 \(b\) 之间放上足够接近的其他量 \(r\) 与 \(s\)，使函数在从 \(r\) 到 \(s\) 的区间上保持连续。这一限制的必要性可以通过如下考虑来体会：级数{[}傅里叶级数{]}的各项都是定积分，而且要回到积分的基本概念。那时将会看到，函数的积分只在函数满足前述条件时才有意义。”

用现代的话说，狄利克雷似乎认为可积性等价于不连续点构成一个稀疏集这一事实；此外，在几行之后，他还指出了那个著名的例子：函数在 x 为有理数时等于 c，在 x 为无理数时取另一个不同的值 d，并断言这个函数“不能被代入”积分。他还宣布要在这一题目上继续工作，但这项工作从未发表，\(^{3}\)而且在此后 25 年间，似乎没有人试图沿这条道路前进，也许是因为在当时看来，考虑这类“病态”函数完全缺乏兴趣；无论如何，当黎曼于 1854 年（{[}259 a{]}，第 227–264 页）重新研究这个问题时（仍然与三角级数相关联\(^{4}\)），他感到有必要为他的工作辩护：“无论我们对力和物质的状态在无限小尺度下随时间与地点变化的方式如何无知，我们仍可确信，狄利克雷的研究所不适用的那些函数并不出现在自然现象中。然而”，他继续写道，“狄利克雷未处理的这些情形值得注意，理由有二。其一，正如狄利克雷本人在其工作末尾所指出的，这个题目与无穷小演算的原理有着极为密切的关系，可以用来给这些原理带来更多的清晰性和确定性。从这个观点看，对它的研究有直接的兴趣。其二，傅里叶级数的应用并不限于物理学中的研究；目前它们在纯粹数学的一个领域即数论中也得到了成功的应用，而在那里，重要的似乎恰恰是那些其三角级数展开尚未被狄利克雷研究过的函数”（{[}259 a{]}，第 237–238 页）。

黎曼的想法是从柯西使之恢复声誉的那个积分近似程序出发，确定一个函数 f 在有界区间 \([a, b]\) 上的“黎曼和”何时趋于一个极限（细分中各区间长度之最大者趋于 0）；他不费力地得到了这个问题如下形式的解：对每个 \(\alpha > 0\)，都存在 \([a, b]\) 的一个分成最大长度足够小的部分区间的分割，使得该分割中 f 的振幅 \(>\alpha\) 的那些区间的长度之和任意小。他还证明，这一条件不仅为连续且逐段单调的函数所满足，而且为可以有处处稠密的不连续点集的函数所满足。\(^{5}\)

黎曼的这篇论文直到他去世后于 1867 年才发表。但这一次，时代对这类研究已较为有利，“黎曼积分”自然而然地汇入了当时正引导人们对“连续统”及实变量函数进行持续研究（魏尔斯特拉斯、杜布瓦–雷蒙、汉克尔、迪尼）、并最终在康托尔那里开出集合论之花的思想潮流。黎曼赋予可积条件的形式，启发人们对函数在一个区间内的不连续点集考虑“测度”的观念；但几乎过了 30 年，才有人给出这一概念的一个富有成效且便于使用的定义。

这个方向上的最初尝试属于施托尔茨、哈纳克和康托尔（1884–1885 年）；前两人为了定义 R 的有界子集 E 的“测度”，考虑包含 E 的集合 \(F \supset E\)，即有限个区间之并，对每个 F 取相应区间长度之和，并把 E 的“测度”说成是这些数的下确界；而康托尔则从一开始就在空间 \(R^{n}\) 中工作，他对有界集 E 以及 \(\rho > 0\) 考虑 E 的邻域 \(V(\rho)\)（由到 E 的距离 \(\leq \rho\) 的点构成），并取 \(V(\rho)\) 的“体积”的下确界。⁶按照这个定义，一个集合的“测度”等于其闭包的“测度”，由此特别可知，两个没有公共点的集合之并的“测度”有可能严格小于这两个集合的“测度”之和。无疑是为了消除后一困难，佩亚诺{[}246 a{]}和若尔当（{[}174 c{]}，第 I 卷，第 28–31 页）在几年之后，对包含在一个方块 I 中的集合 A，在康托尔的“测度” \(\mu(A)\) 之外又引入它的“内测度” \(\mu(I) - \mu(I - A)\)，并把这两个数相等的那些集合 A（现称“可求积的”）称为“可测的”。这时两个没有公共点的可求积集 A、B 之并是可求积的，其“测度”等于 A 与 B 的“测度”之和；但有界开集不一定可求积，包含在一个有界区间中的有理数集亦然，这就使佩亚诺–若尔当的概念失去了大部分意义。

E. 博雷尔{[}32 a{]}的功绩在于能够看出以往各种定义的缺陷，并且看到了如何补救它们。自康托尔以来人们已经知道，R 中的每个开集 U 都是它的可数多个“分支”即两两不交的开区间之并；博雷尔不再试图用有限个区间把 U 围住而从“外部”逼近它，而是依靠上述结果，提议把 U 的各分支长度之和取作 U 的（当 U 有界时的）测度。随后他十分简略地\(^{7}\)描述了可以从开集出发、把可数并和“差” A - B 这两种运算无限迭代而得到的那个集合类（也称“博雷尔”集），并指出对这类集合可以定义一个具有完全可加性这一基本性质的测度：如果序列 \((A_{n})\) 由两两不交的博雷尔集组成，则它们之并（假定有界）的测度等于它们的测度之和。

这个定义将为分析开创一个新纪元：一方面，它与贝尔同时代的工作相配合，成为关于点集分类的一整系列拓扑性质研究的出发点（参见第 165 页）；而尤其重要的是，它将成为勒贝格在 20 世纪初头几年中所创造的积分概念推广的基础。

勒贝格在他的论文{[}196 a{]}中，首先精确化并发展了 E. 博雷尔的简略提示；他仿照佩亚诺–若尔当的方法，把有界集 \(A \subset \mathbf{R}\) 的“外测度”定义为包含 \(A\) 的开集的测度的下确界；然后，若 \(I\) 是包含 \(A\) 的一个区间，则 \(A\) 的“内测度”是 \(I\) 与 \(I - A\) 的外测度之差；这样就得到了“可测集”的概念，它与博雷尔最初的“构造性”定义的差别只在于添上了博雷尔意义下测度为零的集合的一个子集。这个定义立即推广到了空间 \(\mathbf{R}^n\)；而把定积分 \(\int_{a}^{b} f(t) dt\)（函数有界且 \(\geq 0\)）看作由曲线 \(y = f(x)\)、直线 \(x = a, x = b\) 与 \(y = 0\) 所围“面积”的那个旧观念，就这样向所有使上述集合之测度有定义的函数 \(f\) 提供了黎曼积分的一个直接推广。但是

勒贝格的独创性并不那么在于这一推广的思想\(^{8}\)，而在于他发现了关于在这样构想的积分中过渡到极限的基本定理，这个定理在他的工作中表现为测度完全可加性的一个推论；\(^{9}\)他立即认识到它的全部重要性，并把它作为基石，用于他早在 1904 年于著名的《积分论与原函数探求讲义》{[}196 c{]}中所给出的其理论的讲授性阐述。\(^{10}\)

我们在这里无法详细描述勒贝格的结果在无穷小演算古典问题的研究中所带来的不可胜数的进步。勒贝格本人在其论文中就已经把他的理论应用于把长度和面积这些古典概念推广到比通常曲线和曲面更一般的集合；关于这一理论半个世纪以来的可观发展，我们请读者参看 L. 切萨里最近的阐述{[}59{]}。还要提到应用于三角级数的那些成果，它们是勒贝格在其论文完成后几乎立即发展的{[}196 b{]}，并为这门理论开辟了新的视野，对其探索时至今日仍远未结束（参见{[}345{]}）。最后尤其是，空间 \(L^p\) 的定义和菲舍尔–里斯定理（{[}111{]}，{[}260 a and c{]}；参见第 213 页）揭示了积分这个新概念在泛函分析中可能起到的作用；这一作用只能随着这一概念后来的种种推广（我们马上就要谈到）而不断增长。

在此之前，我们要在勒贝格倾注心力最多的那些问题之一上稍作停留，即积分与原函数这两个概念之间的联系。随着黎曼所引入的积分的推广，自然产生的问题是：对连续函数有效的积分与原函数之间的古典对应，在更一般的情形下是否仍然成立。但很容易举出黎曼意义下可积的函数 f 的例子，使得 \(\int_{a}^{x}f(t)dt\) 在某些点没有导数（甚至也没有右导数或左导数）；

反过来，沃尔泰拉早在 1881 年就证明，一个函数 \(F(x)\) 可以在一个区间 \(I\) 中有有界的导数，却在 \(I\) 中（黎曼意义下）不可积。通过一种极为精细的分析（其中关于积分中过渡到极限的定理远远不够用），勒贝格成功地证明：若 \(f\) 在 \([a,b]\) 中（按他的意义）可积，则 \(F(x) = \int_{a}^{x} f(t) dt\) 几乎处处有等于 \(f(x)\) 的导数{[}196 c{]}。反之，若函数 \(g\) 在 \([a,b]\) 中可微且其导数 \(g' = f\) 有界，则 \(f\) 可积，并且我们有公式 \(g(x) - g(a) = \int_{a}^{x} f(t) dt\)。但勒贝格注意到，当 \(g'\) 不有界时问题要复杂得多；这时 \(g'\) 不一定可积，因此第一个问题便是刻画这样的连续函数 \(g\)：其 \(g'\) 几乎处处存在且可积。把自己限制在其诸导数\(^{11}\)之一处处有限的那些函数 \(g\) 上，勒贝格证明了这样的 \(g\) 必为有界变差函数。\(^{12}\)最后他建立了这最后一个结果的逆命题：有界变差函数 \(g\) 几乎处处有导数，且 \(g'\) 可积；但不再必然成立的是

\[
g (x) - g (a) = \int_ {a} ^ {x} g ^ {\prime} (t) d t;\tag{22.2}
\]

这个关系两边之差是一个有界变差函数，它不恒为常数且几乎处处导数为零（一个“奇异”函数）。剩下要刻画的是使关系式 (22.2) 成立的那些有界变差函数 g。勒贝格建立了：这些函数（维塔利称之为“绝对连续”函数，并对它们作了详细研究）正是具有下述性质的函数：g 在开集 U 中的全变差（即 g 在 U 的每个连通分支中的全变差之和）随 U 的测度趋于 0。

我们在下文将会看到，这些结果后来如何以一种弱化了的形式获得远为一般的适用范围。就其最初的形式而言，它们的应用领域仍相当有限，没有超出实变量函数“精细”理论的框架，对这一理论的后续发展我们在这里不作研究；我们只提一下当茹瓦及其追随者和后继者（佩龙、德拉瓦莱–普桑、辛钦、卢津、巴拿赫等）的深刻工作；读者可在 S. 萨克斯的书{[}268{]}中找到详细的阐述。

勒贝格理论所带来的实质性进步之一涉及多重积分。这个概念大约在 18 世纪中叶被引入，而且首先是作为“不定积分”的形式（通过与一元函数积分论的类比，\(\int\int f(x,y)dx dy\) 表示方程 \(\frac{\partial^{2}z}{\partial x\partial y}=f(x,y)\) 的一个解）；但早在 1770 年，欧拉就已经对展布在有界区域（由解析曲线弧围成）上的二重积分有了非常清楚的概念，并且正确地写出了用两个依次进行的简单积分来计算这种积分的公式（{[}108 a{]}，(1)，第 XVII 卷，第 289–315 页）。只要被积函数连续、积分区域也不太复杂，从“黎曼和”出发证明这个公式并不困难；但一旦想处理更一般的情形，黎曼的程序就遇到了严重的困难\((f(x,y)\) 可以在黎曼意义下可积，而当各简单积分按黎曼意义取时，\(\int dx\int f(x,y)dy\) 却没有意义）。当我们转到勒贝格的定义时，这些困难就消失了；勒贝格早在其论文中就证明了：当 \(f(x,y)\) 是有界的“贝尔函数”时，函数 \(y\mapsto f(x,y)\)（对一切 x）和 \(x\mapsto\int f(x,y)dy\) 也如此，并且我们有公式

\[
\int \int f (x, y) d x d y = \int d x \int f (x, y) d y\tag{22.3}
\]

（积分在矩形上取）。稍后，富比尼{[}121{]}为这一结果补充了一个重要的定理：若只假设 \(f\) 可积，则 \(x\) 的那些值（对它们而言 \(y \mapsto f(x, y)\) 不可积）构成一个零测集，这就使公式 (22.3) 得以推广到这种情形。

最后，在 1910 年{[}157 e{]}，勒贝格着手把他关于简单积分之导数的结果推广到多重积分。这样他就被引导着对一个函数 \(f\)（它在 \(\mathbf{R}^n\) 的所有紧子集上可积）联结上集合的函数 \(F(E) = \int_E f(\mathbf{x}) dx\)，它对 \(E\) 取遍 \(\mathbf{R}^n\) 的一切可积子集都有定义，推广了“不定积分”的概念；他借此机会注意到这个函数具有以下两条性质：\(1^\circ\) 它是完全可加的；\(2^\circ\) 它是“绝对连续的”，意思是 \(F(E)\) 随 \(E\) 的测度趋于 0。勒贝格这篇论文的实质部分在于证明这个命题的逆命题。\(^{13}\)但他没有止步于此，而是在同一方向上指出了推广有界变差函数概念的可能性，即考虑可测集的函数 \(F(E)\)，它完全可加，并且 \(\sum_n |F(E_n)|\) 对 \(E\) 分成可测子集 \(E_n\) 的每个可数分割都保持有界。而且，尽管他实际上只限于在 \(\mathbf{R}^n\) 的方块的集合上考虑这样的函数，但十分清楚，只需再走一步就可得到 J. 拉东将在 1913 年定义的一般的测度概念，它把勒贝格积分与斯蒂尔杰斯积分汇集于同一个综合之中，对后一种积分我们现在必须谈一谈。

1894 年，T. 斯蒂尔杰斯以《连分数研究》为题{[}297{]}发表了一篇极具独创性的论文，其中从表面看来十分狭窄的问题出发，以罕见的优雅陈述并解决了在解析函数论和实变量函数论中性质全新的问题。\(^{14}\)为了表示解析函数某个序列的极限，斯蒂尔杰斯（除其他事情外）被引导着在直线上引入正“质量分布”的概念，这个概念在物理科学中久已为人熟悉，但在此之前，它在数学中还从未在限制性假设（一般地，即在每个点存在连续变化的“密度”）之外被考虑过；他指出，给出这样一个分布等价于给出一个递增函数 \(\phi(x)\)，它对 x \textgreater{} 0 给出端点为 0 与 x 的区间中所含的总质量，对 x \textless{} 0 则给出改变符号后的这个质量，而 \(\phi\) 的不连续点对应于“集中在一点上”的质量。\(^{15}\)接着，斯蒂尔杰斯对区间 \([a, b]\) 中的这样一种质量分布构造“黎曼和” \(\sum_{i} f(\xi_{i})(\phi(x_{i+1}) - \phi(x_{i}))\)，并证明当 f 在 \([a, b]\) 中连续时，这些和趋于一个极限，他把它记作 \(\int_{a}^{b} f(x) d\phi(x)\)。由于只需要积分连续函数（甚至是可微函数），斯蒂尔杰斯没有再深入研究这个积分\(^{16}\)，在大约十年间这个概念似乎不曾引起注意。\(^{17}\)但 1909 年，F. 里斯{[}260 d{]}在解决阿达马几年前提出的一个问题（参见第 213 页）时，证明了斯蒂尔杰斯积分 \(f \mapsto \int_{a}^{b} fd\phi\) 是空间 \(\mathcal{C}(I)\)（在 \(I = [a, b] (\mathcal{C}(I))\) 上连续的数值函数的空间，赋有一致收敛拓扑\(^{18}\)）上最一般的连续线性泛函；这个结果的优雅与简洁几乎立即引起了它的几个推广。其中最成功的是 J. 拉东 1913 年的推广{[}256{]}：他把 F. 里斯的思想与勒贝格的思想结合起来，证明了如何用勒贝格的程序、从一个任意的“完全可加的集合函数”（定义在勒贝格测度意义下的可测集上）出发而不是从勒贝格测度出发，来定义一个积分。在这样定义的 \(R^{n}\) 上的“拉东测度”概念中，“有界变差”函数的概念被吸收了：把这样一个函数分解为两个递增函数之差，是测度分解为两个正测度之差的一个特例；同样，“以 \(\mu\) 为基的测度”对应于“绝对连续”函数的概念，而把任意测度分解为一个以 \(\mu\) 为基的测度和一个与 \(\mu\) 互外的测度，对应于有界变差函数分解为绝对连续函数与“奇异”函数之和的勒贝格分解。此外，拉东还证明了，关于 \(\mu\) 的“密度”——当一个测度以 \(\mu\) 为基时——在 \(\mu\) 是以勒贝格测度为基的测度的情形下仍然存在，他所用的是 F. 里斯先前的一个想法（后来被 J. 冯·诺伊曼等人重新提起并普及），即构造测度 \(\mu\) 在一个映射 \(\theta\)（由 \(R^{n}\) 到 R，其选取使 \(\theta(\mu)\) 是 R 上的勒贝格测度）下的象。

拉东的论文发表后几乎立即，弗雷歇就指出，这项工作中的几乎全部结果都可以推广到如下情形：“完全可加的集合函数”不定义在 \(R^{n}\) 的可测集上，而是定义在一个任意集合 E 的某些子集上（这些子集使得可数并与“差”这两种运算仍然给出该函数有定义的集合）。然而，把以 \(\mu\) 为基的测度表成 g. \(\mu\) 的形式这件事，在勒贝格和拉东那里依赖于本质上使用了 \(R^{n}\) 的拓扑的论证（而且我们已经看到，拉东的证明只有当 \(\mu\) 是以勒贝格测度为基的测度时才适用）；直到 1930 年，O. 尼科迪姆{[}235{]}才用直接的论证得到了一般形式的这个定理（几年后 J. 冯·诺伊曼借助空间 \(L^{2}\) 的性质（{[}324 g{]}，第 127–130 页）对它作了显著的简化）。

有了拉东的论文，积分的一般理论就可以认为在其主要轮廓上已经完成；作为其后的实质性收获，只能提到由丹尼尔{[}76 b{]}给出的测度无穷乘积的定义，以及博赫纳 1933 年给出的取值于巴拿赫空间的函数的积分{[}24 a{]}，后者是几年后由盖尔范德{[}125 a{]}、邓福德和佩蒂斯（{[}95{]}与{[}96{]}）所发展的更一般的“弱积分”研究的先声。但还剩下使新理论普及开来、使它成为普遍使用的数学工具的工作，因为在 1910 年前后，大多数数学家在“勒贝格积分”中仍然只看到一件精密度很高、操作细腻、只适用于极端精细和极端抽象的研究的工具。这就是卡拉泰奥多里在一本长期保持经典地位的著作{[}49{]}中所做的工作，这部书以大量独创性的评注充实了拉东的理论。

但也正是在这部书中，一直处于勒贝格关注之首的积分概念（他的论文{[}196 a{]}和他关于这些问题的主要著作{[}196 c{]}的标题就足以表明这一点），第一次把地盘让给了测度概念，而后者对勒贝格来说（一如在他之前对若尔当那样）只是一种辅助的技术手段。观点上的这一变化，在卡拉泰奥多里那里无疑是由于他似乎赋予了“p 维测度”过大的重要性。\(^{19}\)从那时起，论述积分的作者们就在这两种观点之间分成了两派，而且不免卷入了一些引起大量墨水（如果不是大量鲜血的话）之争的论战。一派追随着卡拉泰奥多里；在他们越来越抽象、越来越公理化的阐述中，测度连同它所允许的一切精细加工，不仅扮演着主导角色，而且趋于失去与拓扑结构的联系，而在测度所涉及的大多数问题中，它实际上是与这些结构相连的。另一些阐述则或多或少紧随 W. H. 杨早在 1911 年就已指出、其后由丹尼尔发展的方法，那篇指出该方法的论文可惜很少受人注意{[}339 b{]}。第一种阐述在处理勒贝格积分时，从假定已知的、具有紧支集的连续函数的“柯西积分”出发，依次定义下半连续函数的上积分，然后定义任意数值函数的上积分，由此用纯“泛函”的手段得到可积函数的一个定义，它是仿照勒贝格对集合的定义作出的。丹尼尔在 1918 年（{[}76 a{]}；参见{[}207{]}）把这一阐述（有一些变化）推广到定义在任意集合上的函数；在同一思想脉络中（并与盖尔范德之前谱论中所用的方法密切相关），还必须指出 F. 里斯的论文{[}260 h{]}，它以优雅而简洁的形式给出了序空间理论中在积分理论里起作用的那些结果。

自 1920 年以来源于积分理论的进步，与其说须在读起来多少令人愉快的阐述性著作中去寻找——它们的实质内容不可能有多大的变化——不如说须在应用方面去寻找：概率论（昔日曾是猜测与悖论的托词，自其

被科尔莫哥洛夫公理化{[}183{]}之后，已成为积分论的一个分支，而且是一个有着自己的方法和问题的自主分支）；遍历理论；谱理论和调和分析，因为哈尔发现以他的名字命名的测度以及这一发现所引发的思想运动，已使积分成为群论中最重要的工具之一。
